\documentclass[12pt,a4paper]{article}
\usepackage[utf8]{inputenc}
\usepackage[english]{babel}
\usepackage{geometry}
\usepackage{xcolor}
\usepackage{enumitem}
\usepackage{hyperref}

\geometry{margin=2.5cm}

\title{\textbf{Response to Reviewers' Comments}}
\author{Brandon Lotero}
\date{\today}

\begin{document}

\maketitle

\section*{Acknowledgments}

We express our profound appreciation to the distinguished reviewers
for their thorough and insightful evaluation of this manuscript. Your
expertise and constructive feedback helped us in enhancing the rigor
and clarity of this research. We are grateful
for the time and dedication you have invested in this review process.
Below, we provide comprehensive responses to each of your valuable
comments and recommendations.

\section*{Response to Comments}

\subsection*{Dear MSc. Cristian Alfonso Jiménez Castaño}

We sincerely thank you for all the valuable comments and
recommendations you have provided. Your observations have been
fundamental for improving the quality and clarity of the work.

\subsubsection*{General Observations}

\textbf{Problem 1.} Document cohesion from the problem statement onwards.
\textbf{Response:} The document has been restructured to improve
cohesion between sections, especially in the methodology and
presentation of results.

\textbf{Problem 2.} Lack of bibliographic support in several
subsections of the state of the art.
\textbf{Response:} Appropriate bibliographic references have been
added in all subsections that lacked support.

\textbf{Problem 3.} Absence of justification for the use of
convolutional networks vs. other architectures.
\textbf{Response:} The methodological section has been expanded to
include detailed justification for the choice of U-Net architecture
and comparison with alternatives.

\textbf{Problem 4.} Results and conclusions presented in
separate and disconnected chapters.
\textbf{Response:} The results sections have been unified and
the connection with the research objectives has been improved.

\subsubsection*{Motivation}

\textbf{Problem 5.} Figure 1-1: font size not recognizable at 100\%.
\textbf{Response:} The font size has been increased and
the figure has been resized for better visibility.

\textbf{Problem 6.} Figure 1-2: small letters and page organization.
\textbf{Response:} The figure has been resized and the
layout has been reorganized to avoid blank spaces.

\textbf{Problem 7.} Lines 429 and 440: usability frequency percentage not clear.
\textbf{Response:} The methodology and study universe have been
clarified for the calculation of these percentages.

\subsubsection*{Problem Statement}

\textbf{Problem 8.} Figure 1-3: font size and accommodation.
\textbf{Response:} The font size has been increased and the
distribution on the page has been improved.

\textbf{Problem 9.} Excessive concentration on multiple annotators
without considering other problems.
\textbf{Response:} The section has been expanded to include a
broader view of the state-of-the-art problems.

\textbf{Problem 10.} Research question without clear methodology
to justify U-Net.
\textbf{Response:} The methodological justification
for the choice of U-Net has been developed.

\textbf{Problem 11.} Quick dismissal of other image types without
justification.
\textbf{Response:} Detailed justification has been added for the
focus on histological images.

\subsubsection*{Objectives}

\textbf{Problem 12.} General objective does not begin with verb in infinitive.
\textbf{Response:} The general objective has been reformulated to
begin with a verb in infinitive.

\textbf{Problem 13.} Line 791: interpretability not reflected in
specific objectives.
\textbf{Response:} The specific objectives have been reformulated to
explicitly include interpretability.

\textbf{Problem 14.} Line 792: very strong comparison that requires
extensive experimental framework.
\textbf{Response:} The scope of the comparison has been adjusted to
be realistic with the available experimental framework.

\subsubsection*{Contributions}

\textbf{Problem 15.} Section advances too much without presenting methodologies.
\textbf{Response:} This section has been moved to the end of the document,
after the conclusions.

\subsubsection*{Preliminaries}

\textbf{Problem 16.} Figure 2-2: improve resolution and size.
\textbf{Response:} The resolution has been improved and the figure
has been resized.

\textbf{Problem 17.} Figure 2-3: remove black part of the interface.
\textbf{Response:} The image has been cleaned to show only the
pathological image.

\textbf{Problem 18.} Section 2.2.2: too empty.
\textbf{Response:} It has been completed with visual examples and
detailed explanations.

\textbf{Problem 19.} Section 2.3.1: inconsistency between labels and icons.
\textbf{Response:} Inconsistencies have been corrected and
the connection between paradigms has been improved.

\textbf{Problem 20.} Line 970: clarify that they are machine learning
paradigms in general.
\textbf{Response:} The scope of the mentioned paradigms has been clarified.

\textbf{Problem 21.} Subsection 2.3: too brief.
\textbf{Response:} It has been expanded with additional development and
bibliographic references.

\subsubsection*{Chained Gaussian Processes}

\textbf{Problem 22.} Subsections 3.1.2 and 3.1.3: without bibliographic support.
\textbf{Response:} Appropriate bibliographic references have been added.

\textbf{Problem 23.} Incorrect order of subsections.
\textbf{Response:} It has been reordered to follow the correct logical sequence.

\textbf{Problem 24.} Subsection 3.3: contribution not clear.
\textbf{Response:} The purpose and contribution of
this section have been clarified.

\textbf{Problem 25.} Table 3-1: organize for better visibility.
\textbf{Response:} The table has been reorganized to improve readability.

\subsubsection*{Loss Functions for Multiple Annotators}

\textbf{Problem 26.} Equation 4-10 too long.
\textbf{Response:} It has been divided into intermediate functions for
better readability.

\textbf{Problem 27.} Section 4.1.4: should not be there without experiments.
\textbf{Response:} It has been merged with the previous section for
better coherence.

\subsubsection*{Chained Deep Learning And Tgcess Loss For Image Segmentation}

\textbf{Problem 28.} Numbering begins at zero.
\textbf{Response:} The numbering has been corrected to begin at one.

\textbf{Problem 29.} Figures 5-1 and 5-2 seem identical.
\textbf{Response:} The duplicate figures have been unified.

\textbf{Problem 30.} Lines 1627-1629: it is not clear if the
regularization terms are scalars.
\textbf{Response:} The nature of the regularization terms
in equation 5-1 has been clarified.

\subsection*{Dear Dr. Juan Bernardo Gómez Mendoza}

We sincerely thank you for your positive evaluation and all the valuable
comments and recommendations you have provided. Your observations
have been crucial for improving the quality and completeness of the work.

\subsubsection*{General Evaluation}

\textbf{Problem 31.} Lack of specific tests to illustrate the
effect of each proposed modification.
\textbf{Response:} Ablative experiments have been added to
demonstrate the individual impact of each component of the solution.

\textbf{Problem 32.} Insufficient characterization of the
mathematical implications of the cost function.
\textbf{Response:} The mathematical section has been expanded to
include detailed analysis of the cost function properties.

\textbf{Problem 33.} Lack of clarity in hyperparameter selection.
\textbf{Response:} optuna was used to select the hyperparameters through all
experiments. The scalar experiment in particular was useful for determining the
most meaningful hyperparameters for the model with the help of the Optuna B

\subsubsection*{Form and Structure Observations}

\textbf{Problem 34.} Improvements in the visual presentation of
figures and tables.
\textbf{Response:} All figures and tables have been reviewed and
improved for greater clarity.

\textbf{Problem 35.} Consistency in mathematical notation.
\textbf{Response:} All mathematical notation in the document has been
reviewed and standardized.

\textbf{Problem 36.} Improvements in writing and text clarity.
\textbf{Response:} The writing throughout the document has been
reviewed and improved for greater clarity.

\section*{Conclusion}

We have systematically addressed all the comments and
recommendations provided by both reviewers. The document has
been significantly improved in terms of:

\begin{itemize}
  \item Structural cohesion and logical flow
  \item Bibliographic support and references
  \item Clarity in methodology and justification of decisions
  \item Presentation of results and analysis
  \item Visual quality of figures and tables
  \item Consistency in notation and writing
\end{itemize}

We thank you again for your valuable time and constructive feedback,
which have been fundamental for elevating the quality of
this research work.

\end{document}
